% Assembled from commit c7aa877 of github.com/qedsphere/AAconjecture.
% The text of each part is copied unchanged from notes/distributions/<file>.tex, in this order:
%   threshold_positive/01_oracle_transfer.tex
%   threshold_positive/02_arbitrary_blocks.tex
%   resolution_dynamics/01_theorem_and_setup.tex
%   resolution_dynamics/02_conditional_dependency_bound.tex
%   resolution_dynamics/03_local_posterior_and_simulation.tex
%   resolution_dynamics/04_weighted_parity_maps.tex
%   resolution_structure/04_product_alphabet_completion.tex
%   resolution_dynamics/05_arbitrary_degree.tex
% One edit: the reference to the conventions section in resolution_dynamics/01 points to Section 1.1 of this file.
\documentclass[11pt]{article}
% Compile with LuaLaTeX.
\usepackage[letterpaper,margin=1in,headheight=15pt]{geometry}
\usepackage{amsmath,amssymb,amsthm}
\usepackage{fontspec}
\usepackage{unicode-math}
\setmainfont{TeX Gyre Pagella}
\setsansfont{TeX Gyre Heros}[Scale=MatchLowercase]
\setmonofont{TeX Gyre Cursor}[Scale=MatchLowercase]
\setmathfont{TeX Gyre Pagella Math}
\usepackage{microtype}
\usepackage{xcolor}
\definecolor{ReportBlue}{HTML}{17384B}
\definecolor{LinkBlue}{HTML}{176078}
\usepackage{booktabs,longtable,tabularx,array}
\usepackage{enumitem,fancyvrb,fvextra,upquote,caption}
\usepackage{titlesec,fancyhdr,needspace}
\usepackage{xurl}
\usepackage[colorlinks=true,linkcolor=LinkBlue,urlcolor=LinkBlue,citecolor=LinkBlue,
  bookmarksnumbered=true,pdfencoding=auto]{hyperref}
\usepackage{bookmark}
\setlength{\parindent}{0pt}
\setlength{\parskip}{5pt plus 1pt}
\setlength{\emergencystretch}{2em}
\setlength{\tabcolsep}{5pt}
\setlength{\LTpre}{8pt}
\setlength{\LTpost}{8pt}
\renewcommand{\arraystretch}{1.15}
\setlist{topsep=4pt,itemsep=3pt,parsep=1pt,leftmargin=*}
\setcounter{tocdepth}{2}
\setcounter{secnumdepth}{3}
\clubpenalty=10000
\widowpenalty=10000
\displaywidowpenalty=10000
\allowdisplaybreaks[2]
\titleformat{\section}{\Large\bfseries\color{ReportBlue}}{\thesection}{0.7em}{}
\titleformat{\subsection}{\large\bfseries\color{ReportBlue}}{\thesubsection}{0.6em}{}
\titleformat{\subsubsection}{\normalsize\bfseries}{\thesubsubsection}{0.5em}{}
\titlespacing*{\section}{0pt}{18pt}{8pt}
\titlespacing*{\subsection}{0pt}{13pt}{5pt}
\titlespacing*{\subsubsection}{0pt}{10pt}{4pt}
\captionsetup{font=small,labelfont=bf}
\pagestyle{fancy}
\fancyhf{}
\fancyhead[L]{\small\sffamily Weighted-interaction full-AA proof}
\fancyhead[R]{\small\sffamily November 2026}
\fancyfoot[C]{\small\thepage}
\renewcommand{\headrulewidth}{0.3pt}
\urlstyle{same}
\providecommand{\tightlist}{\setlength{\itemsep}{2pt}\setlength{\parskip}{0pt}}
\providecommand{\pandocbounded}[1]{#1}
\newtheorem{theorem}{Theorem}[section]
\newtheorem{lemma}[theorem]{Lemma}
\newtheorem{proposition}[theorem]{Proposition}
\newtheorem{corollary}[theorem]{Corollary}
\theoremstyle{definition}
\newtheorem{definition}[theorem]{Definition}
\theoremstyle{remark}
\newtheorem{remark}[theorem]{Remark}
\providecommand{\E}{\mathbb E}
\providecommand{\R}{\mathbb R}
\providecommand{\N}{\mathbb N}
\providecommand{\AA}{\mathrm{AA}}
\providecommand{\cAA}{c_{\mathrm{AA}}}
\DeclareMathOperator{\Var}{Var}
\DeclareMathOperator{\Cov}{Cov}
\DeclareMathOperator{\sech}{sech}
\DeclareMathOperator{\tr}{tr}
\DeclareMathOperator{\rank}{rank}
\DeclareMathOperator{\supp}{supp}
\DeclareMathOperator{\diag}{diag}
\pdfstringdefDisableCommands{%
  \def\AA{AA}\def\E{E}\def\Var{Var}\def\Cov{Cov}%
  \def\beta{beta}\def\gamma{gamma}\def\lambda{lambda}%
  \def\mu{mu}\def\rho{rho}\def\le{<=}\def\ge{>=}%
}
% Needed by converted tables.
\usepackage{calc}
\newcounter{none}
\title{Full AA at bounded weighted interaction strength:\\ the assembled proof of Theorem~\ref{thm:resolution-dynamics-arbitrary-degree}}
\author{}
\date{November 2026}
\begin{document}
\maketitle
\thispagestyle{fancy}
\tableofcontents
\clearpage
\section{Conventions, source transfer and block products}

% ---- threshold_positive/01_oracle_transfer.tex ----
\subsection{The two directions a sampling representation must support}
\label{sec:threshold-positive-transfer}

Throughout this section, \emph{full AA} means classical approximation of
the acceptance probability of an arbitrary quantum query algorithm under
the specified input law. We use mean absolute error: a simulator for a
\(T\)-query acceptance function \(Q\) has
\(\mathbb E_\mu|A-Q|\le e\). Substituting \(e=\epsilon\delta\) gives
error at most \(\epsilon\) except on a \(\delta\)-fraction of inputs.
We import the uniform-input simulation consequence of AA, with query
bound \(s_U(T,e)=\operatorname{poly}(T,1/e)\), as in the existing
exchangeable-simulation section. All statements below are deductions
conditional on that imported result, not additional claims about its
provenance. Computation, preprocessing, and the storage of nonuniform
decision trees are free; only queries to the input are counted.

\begin{proposition}[A local posterior transfer criterion]
\label{prop:threshold-local-posterior}
Let \((X,Z)\) be a coupling with \(X\sim\mu\), \(Z\sim\nu\), where
\(\nu\) has a full-AA simulator costing \(s_\nu(T,e)\). Suppose a
deterministic map \(F\) has the following properties.
\begin{enumerate}
\item A coherent query to \(F(Z)_i\) can be implemented with at most
\(a\) queries to \(Z\), uniformly in \(i\).
\item For every \(x\) of positive \(\mu\)-probability and every \(i\),
\(\Pr[F(Z)_i\ne x_i\mid X=x]\le\gamma\).
\item Given oracle access to \(x\), the conditional law \(Z\mid X=x\)
can be supplied consistently on demand: answering any adaptive sequence
of \(s\) source-bit requests uses at most \(b(s)\) queries to \(x\).
\end{enumerate}
Then a \(T\)-query acceptance function has a deterministic simulator
using at most \(b(s_\nu(aT,e/2))\) queries and mean error at most
\[
 e/2+4T\sqrt\gamma.
\]
\end{proposition}

\begin{proof}
Apply simulation under \(\nu\) to the genuine quantum algorithm
\(Q\circ F\), obtaining a source decision tree \(D(Z)\) with average
error \(e/2\). On input \(X=x\), run this tree using the conditional
source oracle from the third hypothesis. Averaging over the coupling
preserves the \(e/2\) bound. The quantum hybrid argument, conditioned
on \(X=x\), gives
\(\mathbb E|Q(X)-Q(F(Z))|\le4T\sqrt\gamma\): at each query the
expected weight on disagreements is at most \(\gamma\), even though
the query weights depend on \(x\). The triangle inequality proves the
claim. The conditional oracle is an ordinary randomized query procedure
using input-independent random coins; fixing those coins by averaging
gives one deterministic simulator of no greater average error.
\end{proof}

The conditional probability in the second hypothesis is essential.
A marginal estimate \(\Pr[F(Z)_i\ne X_i]\le\gamma\) does not by
itself control quantum query weights that adapt to \(X\). Likewise,
an efficient forward sampler alone does not provide the third
hypothesis. The theorem requires sampling the hidden randomness
\emph{conditional on the actual input}, with a query bound. An
arbitrary global posterior computation might require reading all of
\(x\).

This criterion has no source-fanout requirement. One source bit may
affect many output spins, provided that its value can be sampled from
the posterior using few observed spins. That is a substantive
difference from local flip-influence transfer.

% ---- threshold_positive/02_arbitrary_blocks.tex ----
\subsection{Arbitrary independent blocks, including hard constraints}
\label{sec:threshold-positive-blocks}

\begin{theorem}[Block-product full AA without nondegeneracy]
\label{thm:threshold-block-products}
Let the input coordinates be partitioned into blocks of size at most
\(k\), and let \(\mu\) be an arbitrary product of laws on these
blocks. The block laws may depend on dimension and may assign
probability zero to some atoms. For \(T\ge1\) and \(0<e<1\), set
\[
 \eta=\frac{e^2}{64T^2},\qquad
 b=k+\left\lceil\log_2(1/\eta)\right\rceil.
\]
Every \(T\)-query acceptance function admits a deterministic simulator
of mean error at most \(e\), using at most
\[
 S_B(T,e;k):=k\,s_U(2bT,e/2)
\]
queries. In particular, a uniform bound on block size implies full AA,
uniformly over all choices of block laws and temperatures.
\end{theorem}

\begin{proof}
For a block law \(p\) on at most \(2^k\) atoms, put \(N=2^b\).
Give each atom \(u\) the integer count
\(c_u=\lceil(1-\eta)Np(u)\rceil\). Since
\(\sum_uc_u\le(1-\eta)N+2^k\le N\), distribute any remaining
counts arbitrarily. The resulting dyadic law \(p'\) satisfies
\(p'(u)\ge(1-\eta)p(u)\), so it has the form
\[
 p'=(1-\eta)p+\eta r
\]
for another probability law \(r\). A deterministic map of \(b\)
fair bits generates \(p'\) exactly.

Perform this construction independently in each block. Couple
\(X\sim\mu\) to the product law \(Y\) by keeping each entire block
with probability \(1-\eta\), otherwise replacing it by an independent
sample from its \(r\). Thus, conditional on \emph{all} of \(X=x\),
each coordinate changes with probability at most \(\eta\).

Compose the block gadgets with \(Q\). Each virtual input query costs
at most \(2b\) fair-source queries by reading the source block,
computing the requested output bit, and uncomputing. Uniform AA gives
a source tree with error \(e/2\). To supply its source bits on input
\(X\), query the entire physical block when that source block is first
visited, generate its coupled \(Y\)-block, and sample the uniform
source block conditional on that \(Y\)-block. Cache the result.
Different blocks are independent, so this is the correct conditional
source oracle, with at most \(k\) physical queries per first visit.
The transfer proposition gives error
\(e/2+4T\sqrt\eta=e\), followed by derandomization.
\end{proof}

The case \(k=1\) covers products with arbitrarily varying coordinate
biases, including deterministic bits, with no minimum bias assumption.
The case \(k=3\) includes products of laws supported on
\(\{+++ ,---\}\), for which the flip-influence property can fail even
at degree three. Thus the absence of a nondegeneracy assumption is
specific to the full query-simulation target. For a Hamiltonian that
is a sum of disjoint bounded-size components, no finite-temperature
full-AA transition is possible under this theorem; even hard-constraint
limits of the block laws are covered.

\section{Heat-bath histories and the bounded-degree theorems}

% ---- resolution_dynamics/01_theorem_and_setup.tex ----
\subsection{A proved high-temperature regime on bounded-degree graphs}
\label{resolution-dynamics:theorem}

We use the full-AA and uniform-input simulation conventions of
Section~\ref{sec:threshold-positive-transfer}. The following is a deduction from uniform-input
AA and the atom-independent block-product simulation theorem, with all
additional probabilistic and query-access steps proved below.

\begin{theorem}[High-temperature full AA for bounded-degree Ising]
\label{thm:resolution-dynamics-ising}
Fix an integer \(D\ge1\) and a constant \(0\le\theta_*<1\). Consider
arbitrary finite Ising graphs of maximum degree at most \(D\), with
\[
 \mu(x)\propto\exp\left(
     \beta\sum_{ij\in E}J_{ij}x_ix_j+\beta\sum_i h_ix_i\right),
 \qquad
 \Lambda=\max_i\sum_{j:ij\in E}|J_{ij}|.
\]
Assume
\[
 \theta:=D\bigl(1-e^{-4\beta\Lambda}\bigr)\le\theta_*<1.
 \tag{HTQ}
\]
Then every \(T\)-query quantum acceptance function has a deterministic
classical simulator of mean absolute error \(e\), using
\(\operatorname{poly}_{D,\theta_*}(T,1/e)\) queries. The polynomial
is uniform in the number of spins, the graph, coupling signs, and all
finite external fields. The same uniform bound holds after arbitrary
positive-probability pinnings.
\end{theorem}

For \(D>1\), a sufficient temperature condition is
\(\beta\Lambda<\frac14\log(D/(D-1))\), with a fixed strict margin.
The theorem provides an actual open high-temperature interval on every
fixed bounded-degree class. It does not establish the full Dobrushin
range, nor a polynomial whose exponent is independent of \(D\) or the
margin in (HTQ). The case of edgeless graphs follows directly from the
product theorem.

\textbf{Colored heat-bath schedule.}
Fix a proper coloring with \(q\le D+1\) colors. One sweep updates the
colors in order. Within a color, every site independently samples its
spin conditional on the other colors. Because no two sites of that
color interact, this is precisely conditional resampling of that color
class. Its transition kernel \(K_c\) is reversible with respect to
\(\mu\):
\[
 \mu(x)K_c(x,y)=\mu(y)K_c(y,x).
 \tag{DB}
\]
For states agreeing off the color class, both sides are the marginal
weight of the outside spins multiplied by the two conditional weights
of the color-class spins; otherwise both vanish. A full ordered sweep
need not be reversible. Its time reversal is the reverse color order,
which is all that will be used.

\textbf{Finite random maps, with no small-atom assumption.}
For a site \(i\), write
\[
 p_i(\eta)=\mu(X_i=+1\mid X_{N(i)}=\eta)
 =\frac{1+\tanh(\beta h_i+\beta\sum_jJ_{ij}\eta_j)}2.
\]
A heat-bath update can be implemented by the random map
\(M_i(\eta)=+1\) when \(U\le p_i(\eta)\), and \(-1\) otherwise,
for a fresh uniform \(U\in[0,1]\). There are at most
\(2^{|N(i)|}+1\le2^D+1\) different maps: the map changes only as
\(U\) crosses one of the finitely many values \(p_i(\eta)\).
Thus its exact distribution can be encoded by a block of
\(b_0=D+1\) Boolean source bits with an arbitrary block law. Unused
codewords are assigned a constant map so that evaluation is defined on
the entire source cube.

Use independent source blocks for all site-update events in \(L\)
sweeps. Their joint law \(\nu\) is a product of blocks of size at most
\(b_0\); it is covered by Theorem~\ref{thm:threshold-block-products}
uniformly over its atom probabilities. In particular, no discretization
of rare heat-bath probabilities, no bounded-field hypothesis, and no
source precision depending on \(n\) are needed.

Let \(W^0\sim\mu\) independently of these maps, and run the schedule
to obtain the stationary trajectory \(W^0,\ldots,W^{qL}\). Write
\(X=W^{qL}\). Run the same maps from the fixed all-plus initial
configuration, producing \(Y=F_L(Z)\). Then \(X\sim\mu\), while
\(Z\sim\nu\) is exactly the product source law. This coupling,
not a claim of global total-variation mixing, will transfer simulation.

% ---- resolution_dynamics/02_conditional_dependency_bound.tex ----
\subsection{A dependency bound conditional on the complete final input}
\label{resolution-dynamics:conditional-bound}

A random threshold map is \emph{oblivious} when it is constant on
all neighbor assignments. Put
\[
 p_i^-:=\min_\eta p_i(\eta),\qquad
 p_i^+:=\max_\eta p_i(\eta),\qquad
 \rho:=1-e^{-4\beta\Lambda}.
\]
We first need a posterior bound, stronger than the usual unconditional
probability that an update is not oblivious.

\begin{lemma}[Posterior activity bound]
\label{lem:resolution-dynamics-posterior-activity}
Conditional on any complete stationary spin trajectory, the forward
random maps are independent across update events. Each map is
nonoblivious with conditional probability at most \(\rho\).
Consequently, for any fixed set \(A\) of distinct update events and
every final configuration \(x\),
\[
 \Pr[\text{every event in }A\text{ is nonoblivious}\mid X=x]
 \le\rho^{|A|}.
 \tag{PA}
\]
\end{lemma}

\begin{proof}
Given the trajectory, the prior product map law is conditioned on one
separate constraint per event: its map sends the recorded neighbor
configuration to the recorded new spin. These constraints involve
disjoint source blocks, proving conditional independence. The initial
weight \(\mu(W^0)\) is constant under this conditioning.

At an event with recorded neighbors \(\eta\) and new spin \(+1\),
the posterior uniform variable lies in \([0,p_i(\eta)]\). Thus the
nonoblivious probability is
\[
 \frac{p_i(\eta)-p_i^-}{p_i(\eta)}
 \le1-\frac{p_i^-}{p_i^+}.
\]
For new spin \(-1\), it is at most
\(1-(1-p_i^+)/(1-p_i^-)\). These formulas are equally valid for
the finite random-map distribution, by summing its threshold intervals.

The function \(a\mapsto\log[(1+\tanh a)/2]\) has derivative in
\([0,2]\). The local fields in \(p_i\) span an interval of length
at most \(2\beta\Lambda\). The same derivative bound applies to
the logarithm of the probability of \(-1\), with the sign reversed.
Therefore
\[
 \frac{p_i^-}{p_i^+}\ge e^{-4\beta\Lambda},\qquad
 \frac{1-p_i^+}{1-p_i^-}\ge e^{-4\beta\Lambda}.
\]
This proves the uniform posterior activity bound, regardless of the
external field. Multiply it over \(A\) conditional on the trajectory,
then average over trajectories conditional on \(X=x\), obtaining
(PA). No independence conditional on \(X=x\) alone is asserted.
\end{proof}

\begin{lemma}[Conditional coupling error after whole sweeps]
\label{lem:resolution-dynamics-coupling}
For every site \(i\) and every final configuration \(x\),
\[
 \Pr[F_L(Z)_i\ne x_i\mid X=x]\le\theta^L,
 \qquad \theta=D\rho.
 \tag{CE}
\]
\end{lemma}

\begin{proof}
If the coupled final spins disagree at \(i\), their last update map
is nonoblivious, and at least one neighbor disagreed immediately before
that update. Follow such a neighbor to its most recent prior update,
and continue backwards until the initial configuration is reached.
Every visited update is nonoblivious. Update times strictly decrease,
so the events in this path are distinct.

Label each event by its sweep number and site color. Moving backwards
to a neighbor's last update decreases the sweep number by zero or one:
an earlier color was updated in the same sweep, and a later color in
the preceding sweep. A path from a final event in sweep \(L\) to
the initial boundary therefore contains at least \(L\) update events.
This observation is why whole sweeps, rather than an unspecified
global mixing time, are used.

There are at most \(D^{L-1}\) possible length-\(L\) prefixes from
the fixed final event. Their vertices and times are determined by the
schedule and the chosen neighbor sequence, independently of the maps.
By (PA) each such prefix is wholly nonoblivious with probability at
most \(\rho^L\), even conditional on all of \(X=x\). A union bound
gives \(D^{L-1}\rho^L\le\theta^L\). A prefix cannot reach the
initial boundary in fewer than \(L\) events. For an isolated site
the update map is oblivious, so the claimed bound also holds.
\end{proof}

% ---- resolution_dynamics/03_local_posterior_and_simulation.tex ----
\subsection{Local reverse access and completion of the simulation proof}
\label{resolution-dynamics:simulation}

Let \(m=qL\), \(B=D+1\), and take the convenient deterministic
causal-cone bound
\[
 N_L=B^{m+1}.
\]
This intentionally loose bound is independent of \(n\).

\textbf{Forward coherent access.}
To evaluate one final spin of \(F_L\), recursively evaluate the
neighbor spins needed by its most recent update, moving backwards
through at most \(m\) layers. An update has at most \(D\) parents;
an inactive layer merely passes the spin backwards unchanged.
There are at most \(N_L\) update nodes in the full recursion tree.
Reading each map's \(b_0\)-bit code and computing the fixed lookup
tables costs at most \(b_0N_L\) source queries. This can be done on
the entire deterministic cone, without relying on random early
termination. Reversible computation and uncomputation implement a
coherent output query with at most
\[
 a_L=2b_0N_L
\]
source queries, including on unused source codewords.

\textbf{An exact posterior oracle.}
Given oracle access to the actual final input \(x\), sample a
stationary history backwards from \(W^m=x\), using the reverse color
order and independent heat-bath coins. Equation (DB) proves that this
is the exact law of the forward stationary history conditional on its
final state, despite the ordered sweep itself not being reversible.

The reverse history can be evaluated lazily and consistently.
Pre-sample independent reverse coins for every update event. To obtain
\(W_i^t\), recurse towards the final layer. If site \(i\) was not
updated at layer \(t+1\), then \(W_i^t=W_i^{t+1}\). If it was,
the reverse heat-bath coin and the neighbor values at layer \(t+1\)
determine \(W_i^t\). Those neighbors were unchanged by the forward
color update. At the final layer the value is the queried spin
\(x_i\). The recursion is acyclic, has depth at most \(m\), and
uses at most \(N_L\) input queries. Cache its random coins and
computed values across all requests.

After generating this history, independently sample each forward map
conditional on its recorded neighbor assignment and new spin. This
is exactly the disintegration used in
Lemma~\ref{lem:resolution-dynamics-posterior-activity}. To answer a
request for a source bit at event \((t,i)\), only
\(W_i^t\) and \(W_{N(i)}^t\) are needed; the neighbor values equal
their pre-update values. These at most \(D+1=B\) history variables
cost at most \(BN_L\) physical-spin queries. Sample the whole
conditional map block using a fresh preassigned coin and cache it.

This procedure samples \(Z\mid X=x\) exactly. It answers \(s\)
arbitrarily adaptive source-bit requests with at most
\[
 b_L(s)=BN_Ls
\]
queries to \(x\). All reverse-history coins and posterior-map coins
can be fixed in advance independently of the input. Consistency follows
from this single global randomized construction. In particular, there
is no assumption that a freshly explored random cone retains its
unconditioned tail after previous adaptive requests: the bound uses
the entire deterministic cone for every request.

\textbf{Proof of Theorem~\ref{thm:resolution-dynamics-ising}.}
Assume \(T\ge1\) and \(0<e<1\); trivial cases need no queries.
If \(\theta_*=0\), take \(L=1\). Otherwise choose
\[
 L=\max\left\{1,
 \left\lceil\frac{2\log(8T/e)}{\log(1/\theta_*)}\right\rceil
 \right\}.
\]
The product-block source law has simulator cost
\(S_B(\tau,\varepsilon;b_0)\) from
Theorem~\ref{thm:threshold-block-products}. Apply the local posterior
transfer proposition to the coupling above, the map \(F_L\), and
the conditional error bound (CE). The resulting physical simulator
has error at most
\[
 e/2+4T\theta^{L/2}\le e
\]
and uses at most
\[
 B N_L\,S_B(a_LT,e/2;b_0)
 \tag{QC}
\]
queries. Since \(q\le D+1\), \(L=O_{\theta_*}(\log(T/e))\),
and \(N_L=(D+1)^{qL+1}\), (QC) is a polynomial in \(T,1/e\)
for fixed \(D,\theta_*\). Randomness is removed by averaging exactly
as in the local posterior transfer proposition.

Pinning spins deletes them from the remaining interaction graph,
reduces its incident interaction sums, and adds external fields.
The same proof therefore applies with the same \(D,\theta_*\),
which proves the hereditary assertion. \qed

\textbf{Bounded-degree higher-order classical interactions.}
The proof also applies to
\(H(x)=-\sum_A J_A\prod_{i\in A}x_i\). Form the graph joining two
spins when they occur in a common nonzero interaction of size at least
two. Suppose this graph has maximum degree at most \(D\), and set
\(\Lambda=\max_i\sum_{A\ni i,\,|A|\ge2}|J_A|\).
A proper coloring ensures that every interaction contains at most one
spin of each color, so each color update is still independent
conditional resampling. The local field varies by at most
\(2\beta\Lambda\), and all map, posterior, and cone bounds are
unchanged. Arbitrary one-spin fields remain allowed. Thus (HTQ)
proves full AA for this class as well. This extension concerns
classical Gibbs laws, not measured noncommuting quantum Gibbs states.

% ---- resolution_dynamics/04_weighted_parity_maps.tex ----
\subsection{Removing the degree factor from the temperature condition}
\label{resolution-dynamics:weighted-maps}

Threshold maps give a simple proof, but they may inspect every neighbor
whenever they are nonconstant. A different exact update representation
controls the expected number of inspected neighbors by the incident
interaction strength itself.

\begin{theorem}[A degree-independent sufficient temperature]
\label{thm:resolution-dynamics-weighted}
Fix \(D<\infty\) and \(k\ge2\). Let
\[
 H(x)=-\sum_A J_A\prod_{i\in A}x_i
\]
have interactions of size at most \(k\), arbitrary one-spin fields,
and an interaction graph of maximum degree at most \(D\), where two
spins are adjacent when they share a nonzero interaction. Set
\[
 K=\beta\max_i\sum_{A\ni i,\,|A|\ge2}|J_A|.
\]
If \(K\le1/[8(k-1)]\), the Gibbs laws have hereditary full AA with
query bound \(\operatorname{poly}_{D,k}(T,1/e)\), uniformly over all
finite fields and all model parameters satisfying this inequality.
For pair interactions, the sufficient condition is simply
\(\beta\Lambda\le1/8\).
\end{theorem}

The temperature bound does not depend on graph degree, but the query
polynomial still may: this theorem does not cover a degree bound growing
with input dimension.

\textbf{Two elementary norms.}
For a function on a finite Boolean cube, write
\(f=\sum_S\widehat f(S)\chi_S\), and set
\[
 \|f\|_A=\sum_S|\widehat f(S)|,\qquad
 W(f)=\sum_S|S|\,|\widehat f(S)|.
\]
Since \(\chi_S\chi_T=\chi_{S\triangle T}\),
\[
 \|fg\|_A\le\|f\|_A\|g\|_A,\qquad
 W(fg)\le W(f)\|g\|_A+\|f\|_AW(g).
 \tag{WA}
\]
All Fourier coefficients and norms here concern only the at most
\(D\) neighbors of one site.

\begin{lemma}[An exact sparse-support update law]
\label{lem:resolution-dynamics-parity-law}
Let the conditional mean of one spin be
\(m(x)=\tanh(h+u(x))\), where
\(u\) has no constant term,
\(\|u\|_A\le K\le1/8\), and
\(W(u)\le(k-1)K\). There is an exact random map with conditional
mean \(m(x)\) that is either constant or a signed parity of a subset
of the neighbors. If \(S\) is that subset (empty for a constant map),
then, conditional on any neighbor assignment and either resulting spin,
\[
 \mathbb E[|S|\mid\text{neighbors, resulting spin}]
 \le (k-1)\chi(K),\qquad
 \chi(K)=\frac{4e^{2K}K\sec^2K}{(1-\tan K)^2}.
 \tag{PS}
\]
\end{lemma}

\begin{proof}
Write \(r=\tanh h\) and \(v=\tanh u\). The absolute values of
the Taylor coefficients of \(\tanh z\) are the positive Taylor
coefficients of \(\tan z\). This follows from
\(\tanh z=-i\tan(iz)\); positivity of the latter coefficients also
follows recursively from \((\tan z)'=1+(\tan z)^2\). The series
converge absolutely in the norm \(\|\cdot\|_A\) for \(K<\pi/2\).
Using (WA) termwise gives
\[
 \|v\|_A\le t:=\tan K,
 \qquad W(v)\le(k-1)K\sec^2 K.
\]
The addition formula and a convergent geometric series yield
\[
 m=r+(1-r^2)\frac{v}{1+rv},\qquad
 \|m-r\|_A\le\frac{(1-r^2)t}{1-|r|t},\qquad
 W(m)\le\frac{(1-r^2)(k-1)K\sec^2K}{(1-|r|t)^2}.
 \tag{PM}
\]

Write \(m=\sum_S c_S\chi_S\) and
\(A=\sum_{S\ne\varnothing}|c_S|\). Since \(t\le1/3\),
\[
 |c_\varnothing|+A
 \le |r|+\frac{(1-r^2)t}{1-|r|t}\le1.
\]
Thus the following probabilities define an actual random map:
choose \(\operatorname{sign}(c_S)\chi_S\) with probability
\(|c_S|\) for each nonempty \(S\), and choose the constant maps
\(+1\) and \(-1\) with respective probabilities
\((1+c_\varnothing-A)/2\) and
\((1-c_\varnothing-A)/2\). Its pointwise mean is exactly \(m\),
so it is the heat-bath update law.

Every possible output spin, on every neighbor assignment, has
probability at least
\[
 \frac{1-\tanh(|h|+K)}2
 \ge\frac{(1-|r|)e^{-2K}}2.
\]
The unconditioned expected support size is \(W(m)\). Conditioning
on any particular output can increase its expectation by at most the
inverse of that output probability. Divide the last bound in (PM) by
this probability and use \(1+|r|\le2\) and
\(1-|r|t\ge1-t\) to obtain (PS).
\end{proof}

\textbf{Uniform contraction constant.}
For \(0\le K\le1/8\), the elementary inequalities
\[
 e^{2K}\le4/3,\qquad \sec^2K\le16/15,
 \qquad\tan K\le2/15
\]
give \(\chi(K)\le(1280/169)K\). For example,
\(\cos K\ge1-K^2/2\ge127/128\) and \(\sin K\le K\)
prove the tangent bound; \(\cos^2K\ge1-K^2\) proves the secant
bound. Consequently the theorem's assumption gives
\[
 \vartheta:=(k-1)\chi(K)\le160/169<1.
 \tag{PC}
\]

\textbf{Proof of Theorem~\ref{thm:resolution-dynamics-weighted}.}
At every site apply the lemma to its interaction field
\[
 u_i=\beta\sum_{A\ni i,\,|A|\ge2}J_A\chi_{A\setminus\{i\}}.
\]
It satisfies the stated norm bounds. Use independent parity-map
choices in the same proper-color schedule as before. There are at
most \(2^D-1\) nonconstant maps and two constant maps, again encoded
by \(D+1\) bits of an arbitrary source-block law. Conditional on the
complete spin trajectory these map blocks are independent, just as
for the threshold maps.

To trace a disagreement backwards through an event, only neighbors in
its selected subset \(S\) need be considered. Conditional on the
trajectory, the expected number of such parents is at most
\(\vartheta\), uniformly at every event. Expanding all possible
dependency paths, the expected number of surviving paths of length
\(L\) is at most \(\vartheta^L\): along each path event times
strictly decrease, so its map blocks are distinct and independent
under this conditioning; summing over each next parent uses (PS).
Repeated vertices reached by different paths do not affect this
first-moment argument. Any dependence on the initial configuration
after \(L\) whole sweeps requires such a path. Averaging over
trajectories conditional on the full final input therefore gives
\[
 \Pr[F_L(Z)_i\ne X_i\mid X=x]\le\vartheta^L.
\]

All deterministic forward cones and the exact reverse-history
posterior oracle from the previous proof are unchanged. In particular,
sampling a forward parity map conditional on its local neighbors and
new spin uses only those local values, regardless of how small its
prior weight is. Choose \(L\) using the fixed bound \(160/169\),
and apply the same query bound (QC). It is a polynomial for fixed
\(D,k\). Pinning can only reduce interaction sizes, primal degree,
and the sum of incident nonconstant coupling magnitudes; terms that
become fields are allowed freely. Thus the assertion is uniform over
pinnings as well. \qed

\section{Product alphabets with symbol queries}

% ---- resolution_structure/04_product_alphabet_completion.tex ----
\subsection{Completion and full AA for arbitrary product alphabets}
\label{sec:resolution-product-alphabets}

This section removes a source-encoding obstruction. A source coordinate
may take values in an arbitrary finite alphabet, and one source query
returns its entire symbol. Both the quantum and classical algorithms
use this same convention. Alphabet sizes, atom probabilities, and the
number of source coordinates do not enter the bounds.

Let \(\Omega=\prod_{i=1}^n\Omega_i\), let
\(\mu=\bigotimes_i\mu_i\), and delete zero-mass symbols if necessary.
For Boolean \(G\) define its resampling influences by
\[
 I_i^\mu(G)=\Pr[G(X)\ne G(X_{-i},X_i')],
 \qquad I^\mu(G)=\sum_i I_i^\mu(G),
\]
where \(X\sim\mu\) and \(X_i'\sim\mu_i\) independently.
For \(x,y\in\Omega\) write
\(\Delta(x,y)=\{i:x_i\ne y_i\}\).
A fractional certificate for \(G\) at \(x\) is a vector
\(c^x\ge0\) with
\(\sum_{i\in\Delta(x,y)}c_i^x\ge1\) whenever \(G(y)\ne G(x)\).

\begin{lemma}[Product-alphabet certificate inequality]
\label{lem:resolution-alphabet-certificates}
Let \(U_b\subseteq G^{-1}(b)\) have mass \(\delta_b>0\), for
\(b=0,1\). Suppose every \(x\in U_b\) has a fractional certificate
of cost at most \(\gamma_b\). Then
\[
 \max_i I_i^\mu(G)\ge\frac{\delta_0\delta_1}{2\gamma_0\gamma_1}.
 \tag{AC}
\]
The same assertion holds after fixing any set of coordinates.
\end{lemma}

\begin{proof}
The only imported combinatorial ingredient is the finite Boolean-box
counting inequality proved as equation (4) in the proof of Lemma 4 of
\href{https://arxiv.org/abs/2611.12795}
{Huffstutler--Miloschewsky--Podder}.
We state exactly the form used. On a two-point product box \(C\), let
\(J\) be its free coordinates and let \(b_i(C)\) count its
unordered bichromatic \(i\)-edges. Consider some ordered antipodal
pairs \((x,y)\), with \(G(x)=1,G(y)=0\), and sets \(H_y\).
Suppose that replacing any subset of coordinates in \(H_y\cap J\)
of the paired \(x\) by their other box values leaves \(G=1\).
Suppose also that every \(z\in C\) with \(G(z)=1\) agreeing with
\(y\) on \(H_y\) satisfies \(w(\Delta(y,z))\ge1\).
Then the number \(N_C\) of these pairs obeys
\[
 N_C\le\sum_{i\in J}w_i b_i(C).                 \tag{BC}
\]
This is a statement about a finite Boolean cube, with no measure in it.
Relabeling each two-symbol coordinate by a bit preserves every
hypothesis. We now supply the measure-dependent extension.

Fix the stated certificates and define
\[
 w_i=2\gamma_0\,\mathbb E[c_i^X\mid X\in U_1],
 \qquad H_y=\{i:c_i^y>w_i\}\quad(y\in U_0).
\]
Then \(\sum_iw_i\le2\gamma_0\gamma_1\) and
\(w(H_y)\le\gamma_0\). Consequently at least half the conditional
\(\mu\)-mass of \(x\in U_1\) satisfies \(c^x(H_y)<1\).
Mark precisely these pairs \((x,y)\). Their product mass is at least
\(\delta_0\delta_1/2\). A marked pair cannot change the value of
\(G(x)\) by changing coordinates in \(H_y\), by its certificate.
If \(G(z)=1\) and \(z\) agrees with \(y\) on \(H_y\), the
certificate of \(y\) gives
\[
 1\le c^y(\Delta(y,z))\le w(\Delta(y,z)).
\]
Thus (BC) applies to the marked antipodal pairs in every box.

Take all boxes \(C=\prod_i C_i\) with \(|C_i|\in\{1,2\}\),
and assign the weight
\[
 a(C)=\prod_i a_i(C_i),\qquad
 a_i(\{u\})=\mu_i(u)^2,
 \quad a_i(\{u,v\})=\mu_i(u)\mu_i(v)\quad(u\ne v).
\]
Every ordered pair \((x,y)\) is antipodal in exactly one such box:
use \(C_i=\{x_i\}\) if the two symbols agree, and
\(C_i=\{x_i,y_i\}\) otherwise. Its box weight equals
\(\mu(x)\mu(y)\). Hence
\[
 \sum_C a(C)N_C\ge\delta_0\delta_1/2.             \tag{PAIR}
\]
For a fixed unordered \(i\)-edge with outside coordinates \(z\)
and endpoint symbols \(u,v\), summing the weights of its containing
boxes gives
\[
 \mu_i(u)\mu_i(v)
 \prod_{j\ne i}\left(\mu_j(z_j)^2+
               \sum_{a\ne z_j}\mu_j(z_j)\mu_j(a)\right)
 =\mu_i(u)\mu_i(v)\mu_{-i}(z).
\]
It follows that
\[
 \sum_C a(C)b_i(C)=\tfrac12 I_i^\mu(G).            \tag{EDGE}
\]
Multiply (BC) by \(a(C)\), sum, and use (PAIR), (EDGE), and the
bound on \(\sum_iw_i\). The result is
\(\delta_0\delta_1/2\le\gamma_0\gamma_1\max_i I_i^\mu(G)\),
which is (AC). Coordinate conditioning leaves a product measure, so
the final assertion follows from the same proof.
\end{proof}

\begin{theorem}[Completion on arbitrary product alphabets]
\label{thm:resolution-alphabet-completion}
Let \(f:D\to\{0,1\}\), where \(D\subseteq\Omega\), and suppose
there are nonnegative weights \(\lambda^x\), for \(x\in D\), with
\[
 \sum_i\lambda_i^x\le A,
 \qquad
 \sum_{i\in\Delta(x,y)}\min(\lambda_i^x,\lambda_i^y)\ge1
 \quad\text{if }f(x)\ne f(y).
 \tag{ADV}
\]
For every \(\eta>0\) there is a deterministic symbol-query tree of
depth at most \(\lceil128A^3/\eta^2\rceil\) whose unconditional
\(\mu\)-mass of errors on \(D\) is at most \(\eta\).
Constant partial functions need no queries.
\end{theorem}

\begin{proof}
Here are the product-measure details of the completion construction.
For \(a\in f^{-1}(1)\), let
\(d_a(z)=\sum_{i:z_i\ne a_i}\lambda_i^a\), and set
\(g(z)=\min\{1,\min_{a\in f^{-1}(1)}d_a(z)\}\).
Then \(g=0\) on the one-inputs and \(g=1\) on the zero-inputs:
for opposite inputs \(a,b\), (ADV) implies
\(d_a(z)+d_b(z)\ge1\) for every \(z\).
At any \(z\) with \(g(z)<1\), choose a minimizing \(a\).
Changing coordinate \(i\) to any other symbol can increase \(g\)
by at most \(\lambda_i^a\). Resampling symmetry therefore gives
\[
 \sum_i\mathbb E|g(X)-g(X_{-i},X_i')|\le2A.
\]
Integrating thresholds over \([1/4,3/4]\) provides a completion
\(G(z)=\mathbf1[g(z)<t]\) with \(I^\mu(G)\le4A\).
For every original domain point \(x\), the vector \(4\lambda^x\)
is a certificate for \(G\): if \(f(x)=1\), an opposite completion
value forces \(d_x(z)\ge t\); if \(f(x)=0\), a minimizing
one-input at \(z\) and the preceding two-distance inequality force
\(d_x(z)>1-t\). Both lower bounds are at least \(1/4\).

At a decision-tree transcript \(C\), let
\(\delta_b(C)=\mu(f^{-1}(b)\mid C)\) and
\(m(C)=\min_b\delta_b(C)\). The remaining measure is product,
and the restricted certificates have cost at most \(4A\).
Lemma~\ref{lem:resolution-alphabet-certificates} gives a free
coordinate with
\[
 I_i^{\mu\mid C}(G)\ge m(C)^2/(32A^2)
\]
unless \(m(C)=0\). Query that coordinate. The exact potential
identity is
\[
 I^{\mu\mid C}(G)
 -\mathbb E_{X_i\mid C} I^{\mu\mid C,X_i}(G)
 = I_i^{\mu\mid C}(G).
\]
Indeed, the other-coordinate resampling influences average unchanged.
After \(s\) queries output the majority domain value in the current
transcript. Its unconditional domain error is
\(e_s=\mathbb E[m(C_s)]\), and \(e_s\) is nonincreasing.
Telescoping the potential identity and using Jensen's inequality yields
\[
 s e_s^2\le\sum_{t<s}\mathbb E[m(C_t)^2]
 \le32A^2 I^\mu(G)\le128A^3.
\]
The stated depth suffices.
\end{proof}

\begin{lemma}[Quantum adversary weights do not depend on alphabet size]
\label{lem:resolution-symbol-adversary}
A \(t\)-query quantum algorithm deciding a partial Boolean function
with error at most \(1/3\) supplies weights (ADV) with \(A=O(t^2)\),
for arbitrary finite symbol alphabets.
\end{lemma}

\begin{proof}
Purify the algorithm. Let \(p_{r,i}(x)\) be its query probability
at index \(i\) immediately before query \(r\), and set
\(a_i(x)=\sum_{r=1}^t p_{r,i}(x)\); thus \(\sum_i a_i(x)=t\).
For opposite outputs, final states have overlap at most
\(2\sqrt2/3\). A single oracle call changes their overlap in
absolute value by at most
\(2\sum_{i:x_i\ne y_i}\sqrt{p_{r,i}(x)p_{r,i}(y)}\): identical
symbols act identically, and the difference of two unitary symbol
oracles has norm at most two. Telescoping and applying Cauchy--Schwarz
over \(r\) therefore gives
\[
 \sum_{i:x_i\ne y_i}\sqrt{a_i(x)a_i(y)}\ge
 c:=\tfrac12(1-2\sqrt2/3)>0.
\]
A second Cauchy--Schwarz inequality gives
\[
 c^2\le
 \left(\sum_{i:x_i\ne y_i}\min(a_i(x),a_i(y))\right)
 \left(\sum_{i:x_i\ne y_i}\max(a_i(x),a_i(y))\right)
 \le2t\sum_{i:x_i\ne y_i}\min(a_i(x),a_i(y)).
\]
Take \(\lambda_i^x=(2t/c^2)a_i(x)\).
\end{proof}

\begin{corollary}[Product-alphabet full AA with symbol queries]
\label{cor:resolution-product-alphabet-aa}
For every finite product measure on arbitrary finite alphabets,
every \(T\)-query quantum acceptance function \(Q\) has a
deterministic classical symbol-query simulator with mean error at
most \(e\) and \(O(T^6e^{-15})\) queries, for
\(T\ge1\) and \(0<e<1\). The absolute implied constant has no
dependence on alphabet sizes, biases, or number of coordinates.
\end{corollary}

\begin{proof}
Put \(m=\lceil16/e\rceil\), \(a=e/16\), and \(\eta=e/4\).
For each \(j=1,\ldots,m\), define a partial function by thresholding
\(Q\) at \(j/m\), with promise \(|Q-j/m|\ge a\).
Repeating the original quantum algorithm \(O(a^{-2})\) times and
thresholding the empirical acceptance frequency decides this partial
function with error at most \(1/3\). The preceding lemma and theorem
give a deterministic classifier \(D_j\) making
\(O(T^6a^{-12}\eta^{-2})\) symbol queries, whose unconditional
error mass on its promise is at most \(\eta\).

Output \(m^{-1}\sum_jD_j\). For each input, at most \(2am+2\)
thresholds lie within distance \(a\) of its acceptance value.
Rounding the integral of threshold indicators to this grid introduces
error at most \(1/m\). Therefore the mean approximation error is
at most
\[
 2a+3/m+\eta\le9e/16<e.
\]
Running the \(m\) classifiers costs \(O(T^6e^{-15})\) queries.
\end{proof}

The new ingredient relative to the Boolean completion theorem is the
weighted two-point-box summation. The symbol-oracle convention is
essential: this corollary does not claim that a whole large symbol can
be recovered from its bit encoding with one Boolean query. Its role
is as an internal source theorem in a transfer argument where both
forward coherent access and conditional source access are supplied
explicitly.

\section{The weighted-interaction theorem without a degree bound}

% ---- resolution_dynamics/05_arbitrary_degree.tex ----
\subsection{Removing the degree and arity restrictions}
\label{resolution-dynamics:arbitrary-degree}

The product-alphabet theorem permits an entire random update map to be
one source symbol. Truncating the number of explored events then removes
the graph-degree dependence from the preceding proof.

\begin{theorem}[Full AA at bounded weighted interaction strength]
\label{thm:resolution-dynamics-arbitrary-degree}
Let
\[
 H(x)=-\sum_{|A|\ge2}J_A\prod_{i\in A}x_i-\sum_i h_ix_i,
 \qquad
 \omega=\beta\max_i\sum_{\substack{A\ni i\\|A|\ge2}}
                       (|A|-1)|J_A|\le1/8.
\]
For finite parameters and \(\beta\ge0\), every \(T\)-query quantum
acceptance function under \(\mu(x)\propto e^{-\beta H(x)}\) has a
deterministic classical simulator of mean error \(e\), using
\(O(T^{18}e^{-28})\) individual-spin queries. Here \(T\ge1\),
\(0<e<1\), and the implied constant is absolute, independent of
dimension, degree, arity, signs, and fields. The bound is uniform after
arbitrary positive-probability pinnings.
\end{theorem}

This is a sufficient high-temperature theorem, not a sharp threshold
or the entire Dobrushin region. For pair interactions its hypothesis is
\(\beta\max_i\sum_j|J_{ij}|\le1/8\). Computation is free; large
symbols appear only in the internal source reduction.

\begin{proof}
\textbf{Uniform map estimate.}
Write the conditional mean at site \(i\) as
\(m_i(\eta)=\tanh(\beta h_i+u_i(\eta))\). Its unweighted and
weighted Fourier norms satisfy
\[
 \|u_i\|_A\le K_i:=\beta\sum_{A\ni i}|J_A|\le\omega,
 \qquad W(u_i)\le w_i:=\beta\sum_{A\ni i}(|A|-1)|J_A|\le\omega,
\]
where these sums omit one-spin terms. The proof of
Lemma~\ref{lem:resolution-dynamics-parity-law}, replacing
\((k-1)K\) by \(w_i\), gives a random signed-parity-or-constant
heat-bath map \(M_i\). If \(S(M_i)\) is its support and
\(p_i(y\mid\eta)=\Pr[M_i(\eta)=y]\), that proof gives
\[
 \frac{\mathbb E|S(M_i)|}{p_i(y\mid\eta)}
 \le\frac{4e^{2K_i}w_i\sec^2K_i}{(1-\tan K_i)^2}
 \le\frac{1280}{169}w_i
 \le\vartheta:=\frac{160}{169}<1.                 \tag{WR}
\]
The numerator is the \emph{unconditional} expected support size.
Thus both the unconditional and output-conditioned expected support
are at most \(\vartheta\), uniformly over all neighbor assignments.

\textbf{Forward truncation.}
Choose any proper coloring of the interaction graph, possibly with
\(n\) colors. Run \(L\) whole heat-bath sweeps from a stationary
initial state using independent map symbols \(Z\), and call the
final input \(X\). Using the same symbols from all-plus defines
\(F_L(Z)\). Conditional on the complete stationary trajectory, map
symbols remain independent, each constrained only to produce its
recorded output on its recorded neighbors. Their conditional expected
support is bounded by (WR).

A disagreement follows a chain of selected parents backwards through
strictly decreasing event times. Reaching the initial boundary takes
at least \(L\) expansions, since a parent step decreases the sweep
number by at most one. Expanding all paths, conditional independence
along each path and (WR) bound the expected number of paths of length
\(r\) by \(\vartheta^r\). Distinct paths may share events; linearity
of expectation suffices. Thus, conditional on every full final input,
\[
 \Pr[F_L(Z)_i\ne x_i\mid X=x]\le\vartheta^L,
 \qquad
 \mathbb E[N_i\mid X=x]\le(1-\vartheta)^{-1},     \tag{FW}
\]
where \(N_i\) counts all nodes in the unrolled evaluation tree,
including initial-boundary nodes. These inequalities follow first
conditional on the entire trajectory, and then by averaging.

Stop evaluation after \(B\) visited nodes, returning a fixed spin on
overflow. The resulting map \(F_{L,B}\) is defined on the entire
product alphabet. One coherent output query costs at most \(2B\)
symbol queries by computing and uncomputing. Reading a map returns its
whole support; manipulating that support or a long time index is free
computation. Markov's inequality gives
\[
 \Pr[F_{L,B}(Z)_i\ne x_i\mid X=x]
 \le\vartheta^L+[B(1-\vartheta)]^{-1}.             \tag{TF}
\]

\textbf{Posterior access without inspecting all neighbors.}
Given \(x\), construct the stationary history backwards, with
independent reverse parity maps in reverse color order. Each color
kernel is conditional resampling and hence self-adjoint; this is the
exact conditional history law even though a complete ordered sweep
need not be reversible. A history spin is evaluated recursively through
the selected supports of reverse maps, ending at queried spins of
\(x\). Cache maps and computed values.

For a new forward-symbol request at event \((t,i)\), first obtain
its resulting history spin \(y\). Independently propose a map from
its prior law. A constant proposal can be tested without any neighbor;
otherwise obtain only the history spins in its support. Reject if the
output differs from \(y\), and repeat until success. Cache the
accepted forward symbol. Use separate independent proposal tapes for
different forward events, and keep them independent of all reverse
maps.

Conditional on the complete history, this is exact rejection sampling
of the forward-map posterior. If \(\tau\) is the first successful
proposal, then
\[
 \mathbb E\left[\sum_{r=1}^{\tau}|S(M_i^{(r)})|
       \mathrel{\Big|}\text{complete history}\right]
 =\frac{\mathbb E|S(M_i)|}{p_i(y\mid\eta)}
 \le\vartheta.                                   \tag{REJ}
\]
The identity sums the geometric probabilities of previous failures;
it does not assume support size is independent of acceptance. Very
rare outputs may require many constant proposals, which cost no
queries. Finite parameters give positive success probability.

For a source tree with at most \(s\) symbol requests, let \(R\)
count its requested history-spin roots: one resulting spin for each
new source event and all neighbor roots from its proposals. Conditional
on the complete history, a newly requested event's proposal tape is
independent of the earlier source transcript. Therefore (REJ) also
applies to adaptively chosen events, giving
\[
 \mathbb E[R\mid X=x]\le s(1+\vartheta).          \tag{ROOT}
\]

Now use the filtration of the \emph{actual lazy computation}, not
conditioning on its unexposed history. A fresh reverse map is still
independent of exposed information: every returned source symbol was
determined from exposed history values and independent proposal coins.
It creates at most \(\vartheta\) children in conditional expectation.
Let \(N\) count processed history-request nodes, including cache hits
and final-boundary queries. Queue accounting, stopped after \(M\)
nodes, yields
\[
 (1-\vartheta)\mathbb E[\min(N,M)\mid x]
 \le\mathbb E[R\mid x].
\]
Indeed roots contribute \(R\) nodes, and each inspected fresh reverse
map contributes its selected children; there are at most
\(\min(N,M)\) fresh inspections. Monotone convergence gives
\[
 \mathbb E[N\mid X=x]
 \le\frac{s(1+\vartheta)}{1-\vartheta}.           \tag{POST}
\]
Every physical query is a boundary-node request. The complete-history
conditioning in (REJ) bounds the number of requested roots; the exposed
information filtration bounds reverse exploration. Independence under
one conditioning has not been transferred to the other.

\textbf{Error and query bounds.}
Set \(\delta=e^2/(288T^2)\) and choose
\[
 L=\left\lceil\frac{\log(1/\delta)}{\log(1/\vartheta)}\right\rceil,
 \qquad B=\left\lceil[\delta(1-\vartheta)]^{-1}\right\rceil.
\]
Then (TF) is at most \(2\delta\), and the quantum hybrid error is
\(4T\sqrt{2\delta}=e/3\). Apply
Corollary~\ref{cor:resolution-product-alphabet-aa} to
\(Q\circ F_{L,B}\) at error \(e/3\). Its source depth is
\[
 s=O((BT)^6e^{-15})=O(T^{18}e^{-27}).
\]
Run it with the posterior oracle, but abort and return zero if history
work exceeds
\(M=\lceil3s(1+\vartheta)/(e(1-\vartheta))\rceil\).
By (POST) the abort probability is at most \(e/3\), conditional on
every input. The uncapped procedure supplies the exact source
posterior, so source simulation, hybrid error, and abort error sum to
at most \(e\). Physical queries are bounded by
\(M=O(T^{18}e^{-28})\).

All coins can be fixed in advance independently of \(x\). Averaging
gives a deterministic simulator; for each finite input space,
almost-sure termination can simultaneously be required at every input.
Finally, pinning reduces interaction support sizes; merging residual
terms obeys the triangle inequality. It therefore cannot increase
\(\sum_{A\ni i}(|A|-1)|J_A|\). The same absolute bound applies
under every pinning.
\end{proof}

\end{document}
